\documentclass[11pt]{article}
\usepackage[a4paper,margin=1in]{geometry}
\usepackage{fontspec}
\usepackage[boldfont=false]{xeCJK}
\setCJKmainfont{FandolSong-Regular.otf}
\usepackage{amsmath}
\usepackage{amssymb}
\usepackage{array}
\usepackage{booktabs}
\usepackage{graphicx}
\usepackage{adjustbox}
\usepackage{enumitem}
\setlist[itemize]{topsep=2pt,partopsep=0pt,parsep=0pt,itemsep=2pt,leftmargin=1.4em}
\setlist[enumerate]{topsep=2pt,partopsep=0pt,parsep=0pt,itemsep=2pt,leftmargin=1.6em}
\usepackage{parskip}
\usepackage{fvextra}
\fvset{breaklines=true,breakanywhere=true,fontsize=\small}
\setlength{\parindent}{0pt}

\begin{document}

\begin{center}
  {\LARGE\bfseries The macroeconomy (A Level)}\\[0.35em]
  {\large A Level Economics}
\end{center}
\vspace{0.6em}

This topic studies the wider economy at A Level depth: how spending multiplies, how growth links to the environment, the jobs–prices link, and money.

\section*{The circular flow and the multiplier}

The \textbf{circular flow of income} 收入循环流 shows money passing between households and firms. \textbf{injections} 注入 add money to the flow (investment, government spending, exports); \textbf{leakages} 漏出 take it out (saving, taxation, imports).

\subsection*{The multiplier}

An injection raises income by \textbf{more} than the first amount. This is the \textbf{multiplier} 乘数. When a firm builds a factory, the builders earn income, and they spend part of it, which becomes someone else's income, and so on.

How much is spent again depends on the \textbf{marginal propensity to consume} 边际消费倾向 (MPC) — the fraction of extra income that people spend on home goods. The rest leaks out. In a simple model:

\[
\text{multiplier} = \frac{1}{1 - MPC}
\]

A bigger MPC (smaller leakages) gives a bigger multiplier. The main leakages are measured by the \textbf{marginal propensity to save} 边际储蓄倾向 (MPS) and the \textbf{marginal propensity to import} 边际进口倾向 (MPM).

\subsection*{The accelerator}

The \textbf{accelerator} 加速数 says that \textbf{investment} 投资 depends on the \textbf{rate of change} of output. When demand is rising quickly, firms need lots of new capital, so investment jumps. When demand merely stays high, investment falls back.

\subsection*{What drives consumption, saving and investment}

\begin{itemize}
\item \textbf{consumption} 消费 and \textbf{saving} 储蓄 depend on \textbf{disposable income} 可支配收入 (income after tax), wealth, interest rates, and confidence.

\item investment depends on interest rates, expected profit, technology, and business confidence.

\end{itemize}

\section*{Economic growth and sustainability}

\textbf{economic growth} 经济增长 is a rise in real output.

\begin{itemize}
\item \textbf{actual growth} 实际增长 is the real rise in output that happens.

\item \textbf{potential growth} 潜在增长 is the rise in the economy's capacity.

\end{itemize}

The \textbf{output gap} 产出缺口 is the difference between actual output and potential output. A \textbf{positive} gap (actual above potential) tends to cause inflation; a \textbf{negative} gap (actual below potential) means spare capacity and unemployment.

\subsection*{Sustainability}

\textbf{sustainable growth} 可持续增长 is growth that meets today's needs without harming the ability of future generations to meet theirs. Fast growth can use up non-renewable resources and cause pollution, which threatens \textbf{sustainability} 可持续性. So a country must weigh growth today against the environment and resources of tomorrow.

\section*{Employment and unemployment}

\textbf{unemployment} 失业 is measured in two ways: the \textbf{claimant count} 申领救济人数 (people claiming benefits) and the wider \textbf{labour force survey} 劳动力调查 (a sample asked about work).

\subsection*{Types}

\begin{itemize}
\item \textbf{frictional unemployment} 摩擦性失业 — short spells between jobs.

\item \textbf{structural unemployment} 结构性失业 — skills no longer match the jobs, as industries decline.

\item \textbf{cyclical unemployment} 周期性失业 — caused by low demand in a recession.

\item \textbf{seasonal unemployment} 季节性失业 — work only at certain times of year.

\end{itemize}

\subsection*{The natural rate and the Phillips curve}

The \textbf{natural rate of unemployment} 自然失业率 is the unemployment left when the labour market is in balance (mainly frictional and structural). It cannot be removed just by raising demand.

The \textbf{Phillips curve} 菲利普斯曲线 shows the link between unemployment and inflation:

\begin{itemize}
\item in the \textbf{short run}, there is a trade-off: lower unemployment comes with higher inflation, because strong demand pushes up wages and prices.

\item in the \textbf{long run}, the curve is vertical at the natural rate. Once people expect inflation, there is no lasting trade-off — pushing unemployment below the natural rate only speeds up inflation.

\end{itemize}

\section*{Money and banking}

\subsection*{What money does}

\textbf{money} 货币 is anything widely accepted in payment. It has four jobs:

\begin{itemize}
\item a \textbf{medium of exchange} 交换媒介 — used to buy and sell, avoiding barter.

\item a \textbf{store of value} 价值储藏 — it keeps its worth over time, so you can save it.

\item a \textbf{unit of account} 记账单位 — a common way to measure and compare prices.

\item a \textbf{standard of deferred payment} 延期支付标准 — it lets people borrow and repay later.

\end{itemize}

Good money is durable, easy to carry, easy to divide, hard to fake, and limited in supply.

\subsection*{Banks}

\begin{itemize}
\item the \textbf{central bank} 中央银行 issues notes, runs monetary policy, holds the country's reserves, and acts as \textbf{lender of last resort} 最后贷款人 to banks in trouble.

\item \textbf{commercial banks} 商业银行 take deposits and make loans. By lending out most of what is deposited, the banking system creates money.

\end{itemize}

\subsection*{The demand for money and interest rates}

People hold money rather than other assets for its \textbf{liquidity} 流动性 (it can be spent at once). The \textbf{demand for money} 货币需求 — also called \textbf{liquidity preference} 流动性偏好 — comes from the need to make everyday payments, to keep a safety buffer, and to be ready to buy assets.

The \textbf{interest rate} 利率 is the price of borrowing money. It is set where the demand for money meets the \textbf{money supply} 货币供给. A larger money supply tends to lower the interest rate.

\subsection*{Moral hazard}

\textbf{moral hazard} 道德风险 is the danger that banks take bigger risks because they expect to be rescued if things go wrong. If a bank believes the government will bail it out, it may lend recklessly — knowing it keeps the gains but passes losses to others.


\end{document}
